

\begin{frame}{Niveaux d'isolation SQL}

La norme SQL définit 4 niveaux.
Un ``vrai''  SGBD en propose au moins deux

\vfill

\begin{redbullet}
  \item \alert{Isolation totale} (\texttt{serializable}). 
  
     Garantit l'ACIDité, mais augmente les verrouillages, parfois l'interblocage
   
  \item \alert{Isolation partielle} (\texttt{read committed} 
     ou \texttt{repeatable read}) 
     
     Permet des anomalies (rares), évite les interblocages.
\end{redbullet}
 
 \vfill
 
 Le choix est une question de compromis, sachant que l'isolation,
 par défaut, est partielle

\end{frame}

\begin{frame}{L'isolation partielle}

 Le mode \texttt{read committed}

\vfill
   \alert{Retenir}: une requête accède à l'état de la base \alert{au moment où la
requête est exécutée}.
  
\vfill

Le mode \texttt{repeatable read}

\vfill

\alert{Retenir}: une requête accède à l'état de la base \alert{au moment où la
transaction a débuté}.

\vfill

\alert{Anomalie possible}

\begin{redbullet}
  \item Deux transactions lisent puis modifient le même nuplet 
   \item La seconde modification annule la première (qui est donc perdue)
\end{redbullet}
\end{frame}


\begin{frame}[fragile]{Isolation totale}

C'est le mode \texttt{serializable}. Il
garantit  la cohérence de la base.

\vfill

\begin{minted}[frame=leftline,
               baselinestretch=.8,
               fontsize=\small,
               framesep=2mm]{sql}
set transaction isolation level serializable;
\end{minted}

\vfill

\alert{Interblocage possible}

\vfill


\begin{redbullet}
  \item Deux transactions lisent le même nuplet 
  
   \item Puis toutes deux veulent modifier ce nuplet, et s'interbloquent

\end{redbullet}
\end{frame}

\begin{frame}{Ce qu'on doit savoir }

Avec un niveau d'isolation partielle, l'exécution est plus fluide,  
les anomalies sont possibles.

\vfill

En isolation totale, la cohérence est assurée, mais l'exécution risque de rencontrer des blocages,
\vfill

\alert{Le niveau d'isolation par défaut n'est jamais le plus strict}. Car

\begin{redbullet}
  \item inutile dans la grande majorité des cas;
  \item provoque rejets et blocages difficiles à expliquer à l'utilisateur.  
\end{redbullet}

\vfill

\alert{Méfiance}\,: en isolation partielle, un programme testé 1\,000 fois
peut malgré tout engendrer des anomalies transactionnelles

\end{frame}
